\chapter{Maryam Mirzakhani}\label{chap:MaryamBio}

\section{A biography}

Maryam Mirzakhani’s biography tells the story of an extraordinary journey: from a young girl in Tehran who dreamed of becoming a writer, to a world-renowned mathematician whose name is firmly established in the history of mathematics. Her path was shaped by determination, and an intense pursuit of truth and beauty. Though cut short by illness, her story continues to inspire mathematicians and non-mathematicians. She remains not only a pioneer in her field but also a symbol of what can be achieved when talent and perseverance are nurtured by opportunity and courage.

\subsection{Early Life}
Maryam Mirzakhani was born on May 3rd, 1977, in Tehran, Iran, to a supportive and intellectually engaged family. Her father, Ahmad Mirzakhani, was an engineer, while her mother, Zahra Haghighi, dedicated herself to raising their four children. Growing up in the years following the Iranian Revolution and the Iran–Iraq War, Mirzakhani later reflected on her good fortune in reaching adolescence during a period of relative stability, as the country worked to rebuild its educational and cultural institutions (\cite{gibbons,series}).

As a child, Mirzakhani did not initially appear destined for a career in mathematics. She was an enthusiastic and avid reader with aspirations of becoming a writer, often spending countless hours absorbed in novels and biographies. She admired the remarkable lives of figures like Marie Curie and Helen Keller, whose stories inspired her with the belief that one could overcome challenges and make a significant impact on the world. Her vivid imagination even led her to create stories about ambitious heroines who traveled extensively and achieved great goals (\cite{gibbons}).

Her introduction to mathematics was partly inspired by her older brother, who one day shared the famous story of Carl Friedrich Gauss summing the numbers from 1 to 100. Although she was not able to solve the problem herself at the time, she was captivated by the idea that such funny shortcuts existed in mathematics. Such interest intensified when she began attending Farzanegan School, a prestigious institution under Iran’s National Organization for Development of Exceptional Talents (NODET). At first, mathematics was not her strongest subject, and a teacher even suggested she lacked aptitude. However, the following year, she had a more encouraging teacher whose belief in her potential reignited her interest (\cite{gibbons}).

Alongside her close friend Roya Beheshti, Mirzakhani successfully convinced their school principal to allow them to train for the International Mathematical Olympiad (IMO). Prior to this, no Iranian girl had ever participated in the competition. In 1994, Mirzakhani won a gold medal at the IMO in Hong Kong, and the following year, she again earned gold in Toronto, with a perfect score. With this accomplishment, she became the first Iranian student, regardless of gender, to obtain such a distinction (\cite{gibbons,series}).

\subsection{University Studies and Academic Formation}
Following her success in international competitions, Mirzakhani enrolled at Sharif University of Technology in Tehran, one of the country’s best institutions for science and engineering (\cite{series}). There, she thrived in a stimulating intellectual environment, enriched by problem-solving sessions and informal reading groups that contributed to her development. As an undergraduate, she published several research papers in combinatorics and graph theory, showing her early originality and determination (\cite{gibbons}).

In 1999, she earned a Bachelor of Science degree in Mathematics. That same year, she left Iran to pursue graduate studies in the United States. She enrolled at Harvard University, where she had the privilege of studying under Curtis T. McMullen, a Fields Medalist (\cite{gibbons,series}). Adapting to a new country and language posed challenges, yet she impressed her professors and peers with her curiosity, frequently posing questions in English while taking notes in Farsi.

Her doctoral years combined by both challenges and achievements. At first, she was not well-versed in certain advanced topics, but her determination and curiosity enabled her to overcome these difficulties. By the time she earned her Ph.D. in 2004, she had gained international fame for the depth and originality of her research (\cite{series}).

\subsection{Professional Career}
After completing her doctorate, Mirzakhani became a research fellow at the Clay Mathematics Institute while also serving as an assistant professor at Princeton University. In 2008, she joined the faculty at Stanford University, where she rapidly ascended the academic ranks and became a full professor by the age of thirty-one: an extraordinary accomplishment in a field known for its competitiveness and rigorous standards (\cite{gibbons,series}).

At Stanford, she distinguished herself not only as an exceptional researcher but also as a devoted educator and mentor. Students and colleagues noted her humility, warmth, and generosity. Although she was reserved in public settings, she became animated when discussing her ideas, often filling pages with intricate sketches as she developed her thoughts (\cite{zagrebnov}).

Her career was recognized by numerous honors. In 2009, she was honored with the American Mathematical Society’s Blumenthal Award, followed by the Ruth Lyttle Satter Prize in 2013. The pinnacle of her recognition came in 2014 when she received the Fields Medal, the most prestigious award in mathematics. This distinguished her as both the first woman and the first Iranian to receive this award since its establishment in 1936 (\cite{zagrebnov}).

Despite such recognitions, Mirzakhani maintained a sense of modesty. She often highlighted the collaborative and playful aspects of mathematics, where persistence and imagination are vital (\cite{gibbons}).

\subsection{Personal Life}
Beyond her professional achievements, Mirzakhani cherished her roles as a wife and mother. In 2005, she married Jan Vondrák, a Czech computer scientist, and in 2011, they welcomed their daughter, Anahita. Friends and colleagues fondly recalled her dedication to her family and her ability to balance the demands of the academic world with her personal commitments (\cite{zagrebnov}).

In 2013, Mirzakhani was diagnosed with breast cancer. Despite her illness, she continued to teach, conduct research, and actively contribute to the mathematical community with remarkable resilience. For several years, she fought the disease while maintaining a spirit of optimism. Tragically, on July 14, 2017, she passed away at the age of forty. Her death was met with profound sorrow across the globe (\cite{zagrebnov,series}).

\subsection{Legacy}
Maryam Mirzakhani's legacy extends far beyond her research contributions. She became a symbol of perseverance and inspiration for women in mathematics and science around the globe (\cite{zagrebnov}). In Iran, her achievements were celebrated as a source of national pride, while internationally, she was recognized as a pioneer who broke through barriers.

Her life is as a powerful reminder that mathematics is not only an abstract discipline but a human endeavor shaped by creativity, imagination, and resilience. She once noted, “The beauty of mathematics only shows itself to more patient followers.” Through her patience and vision, she revealed that beauty to the world (\cite{series}).

Her memory lives on through the many students she mentored, the colleagues who respected her, and the countless young individuals who view her as a model of what can be accomplished through passion and determination. Academic institutions and international societies have established prizes, lectures, and memorials in her honor, ensuring that her influence continues to thrive.

\section{Mathematical Contributions}

The mathematical work of Maryam Mirzakhani constitutes a systematic synthesis of hyperbolic geometry, complex analysis, topology, and dynamical systems. Her research is characterized by a geometric approach that reformulates complex algebraic and counting problems into global questions regarding the topology and geometry of moduli spaces.\footnote{Mirzakhani completed her Ph.D. at Harvard University in 2004 under the supervision of Curtis McMullen. Her early research culminated in her 2007 papers on Weil-Petersson volumes, establishing the foundation for her later work on the moduli spaces of bordered Riemann surfaces \cite{Mirzakhani2007a, Mirzakhani2007b}.}

A unifying theme of her research is the translation of local geometric behaviors—such as the trajectory of a single curve—into global topological invariants. This methodology revealed rigorous connections between pure geometry, topological classification, and mathematical physics.

\subsection{Hyperbolic Geodesics: Beyond the Prime Geodesic Theorem}

The study of geodesics on hyperbolic surfaces is a central topic in differential geometry and ergodic theory. By the Uniformization Theorem, any closed topological surface $S$ of genus $g \geq 2$ admits a hyperbolic metric of constant negative curvature $\kappa = -1$. On such surfaces, the distribution of closed geodesics exhibits chaotic behavior characterized by exponential growth.

It is a well-established result in dynamical systems that the number of all closed geodesics on a hyperbolic surface grows exponentially with respect to their length $L$.\footnote{The Prime Geodesic Theorem states that the number of closed geodesics $\pi(L)$ of length at most $L$ grows asymptotically as $\pi(L) \sim e^{L}/L$. This result is traditionally derived using the Selberg trace formula, establishing a deep analogy with the Prime Number Theorem in number theory.} Mirzakhani addressed the more constrained problem of counting \textit{simple} closed geodesics, defined as those without self-intersections.

Counting simple geodesics presents a distinct topological challenge: as the length $L$ increases, a generic geodesic tends to become equidistributed on the surface, making the "simple" condition statistically rare. Mirzakhani proved that the number of simple closed geodesics $s_X(L)$ grows polynomially rather than exponentially. Specifically, she established the following asymptotic limit:
\begin{equation}
    \lim_{L\to\infty} \frac{s_X(L)}{L^{6g-6}} = c_X
\end{equation}
where $6g-6$ corresponds to the dimension of the relevant Teichmüller space $\mathcal{T}_g$.\footnote{The constant $c_X$ depends on the hyperbolic structure of $X$. It is expressed as $c_X = \int_{\mathcal{M}_g} f_X \omega_{WP}^{3g-3}$, where $f_X$ is the "frequency function" introduced by Mirzakhani to relate the length of simple closed curves to the Weil-Petersson geometry \cite{Mirzakhani2007a}.}

\subsection{Pants Decomposition and Weil--Petersson Volumes}

To analyze the structure of the moduli space $\mathcal{M}_{g,n}$ (the parameter space of all possible hyperbolic structures) and compute the constants $c_X$, Mirzakhani utilized the method of \textit{pants decomposition}. Any hyperbolic surface $X$ of genus $g$ with $n$ boundary components can be partitioned along $3g-3+n$ disjoint simple closed geodesics into $2g-2+n$ \textit{pairs of pants}. 

By parameterizing the surface through these blocks, Mirzakhani established a correspondence between the discrete problem of counting geodesics and the continuous evaluation of integrals over $\mathcal{M}_{g,n}$ equipped with the Weil-Petersson metric. The volume $V_{g,n}(\mathbf{L})$ is obtained by integrating the top power of the associated symplectic form $\omega_{WP}$ over the moduli space:
\begin{equation}
V_{g,n}(\mathbf{L}) = \frac{1}{d!} \int_{\mathcal{M}_{g,n}} \omega_{WP}^d
\end{equation}
where $d = 3g-3+n$ is the complex dimension of $\mathcal{M}_{g,n}$.\footnote{For a complete derivation of the volume form and its integration over the moduli space of bordered Riemann surfaces, see \cite{Mirzakhani2007a}.} A key analytic tool in this framework is the generalization of the \textit{McShane identity} to surfaces with geodesic boundaries $(\ell_1, \dots, \ell_n)$ \cite{Mirzakhani2007a}.

This decomposition leads to a recursive formula for the Weil–Petersson volumes. By removing a single pair of pants, the volume is expressed through the following Mirzakhani recursion:

\begin{equation}
\begin{aligned}
V_{g,n}(\mathbf{L}) &= \sum_{i=2}^{n} \int_{0}^{\infty} V_{g,n-1}(\mathbf{L}_{\hat{1,i}}, \ell) \cdot \ell \, K(\ell_1, \ell_i, \ell) \, d\ell \\
&+ \int_{0}^{\infty} \int_{0}^{\infty} V_{g-1,n+1}(\mathbf{L}_{\hat{1}}, \ell, \ell') \cdot \ell \ell' \, K(\ell_1, \ell, \ell') \, d\ell d\ell' \\
&+ \sum_{\substack{g_1+g_2=g \\ I \sqcup J = \{2,\dots,n\}}} \int_{0}^{\infty} \int_{0}^{\infty} V_{g_1,|I|+1}(\mathbf{L}_I, \ell) V_{g_2,|J|+1}(\mathbf{L}_J, \ell') \cdot \ell \ell' \, K(\ell_1, \ell, \ell') \, d\ell d\ell'
\end{aligned}
\end{equation}

where $\mathbf{L} = (\ell_1, \dots, \ell_n)$.\footnote{The first and second terms correspond to the non-separating case (reducing genus or boundary count), while the third term corresponds to the separating case where the surface splits into components of lower complexity \cite{Mirzakhani2007a}.}

\subsection{A Geometric Proof of Witten's Conjecture}

The recursive framework for Weil--Petersson volumes provided a new approach to Witten's 1991 conjecture regarding the intersection numbers of tautological classes on the Deligne--Mumford compactification $\overline{\mathcal{M}}_{g,n}$. Mirzakhani established that the volumes $V_{g,n}(L)$ are polynomials in $L_i^2$, whose coefficients are related to the intersection numbers of $\psi$-classes:
\begin{equation}
V_{g,n}(L) = \sum_{|\alpha| \leq d} \frac{\langle \tau_{\alpha_1} \dots \tau_{\alpha_n} \rangle}{\alpha! 2^{|\alpha|}} L_1^{2\alpha_1} \dots L_n^{2\alpha_n}
\end{equation}

Mirzakhani demonstrated that her recursive volume formula is equivalent to the Virasoro constraints, which imply the Korteweg-de Vries (KdV) hierarchy.\footnote{Specifically, after applying a multi-dimensional Laplace transform $\mathcal{L}$ to $V_{g,n}$, one obtains a generating function $F(t_0, t_1, \dots)$ satisfying the Virasoro constraints, which are equivalent to the KdV recursion \cite{Mirzakhani2007b}.}

\subsection{Dynamics on Moduli Spaces: The "Magic Wand" Theorem}

In her later research, Mirzakhani shifted to the dynamics of translation surfaces—geometric objects formed by identifying opposite sides of polygons in the plane. In collaboration with Alex Eskin and Amir Mohammadi, she addressed the classification of orbit closures under the $SL(2,\mathbb{R})$ action on the moduli space of Abelian differentials. 

\begin{theorem}[Eskin--Mirzakhani--Mohammadi, 2015]
Every $SL(2,\mathbb{R})$-orbit closure in the moduli space of Abelian differentials is an affine invariant submanifold.
\end{theorem}

This result implies that orbit closures are locally defined by linear equations with real coefficients in period coordinates.\footnote{This theorem represents a deep non-homogeneous analogue of Ratner’s theorems for unipotent flows on homogeneous spaces \cite{EskinMirzakhaniMohammadi2015}.} 

\subsection{Large Genus Asymptotics}

Mirzakhani established precise asymptotic bounds for volumes as the genus $g \to \infty$, a critical limit for connections with statistical mechanics and quantum gravity:
\begin{equation}
    V_g \sim C \cdot (2g)! \cdot (2\pi^2)^{-g} \cdot g^{-1/2}
\end{equation}
where $C$ is a universal constant.\footnote{See \cite{Mirzakhani2013} for the proof of this $g^{2g}$ growth rate, which characterizes the statistical properties of "typical" hyperbolic surfaces of high genus.}

\subsection*{Summary and Methodological Implications}

The primary milestones of Maryam Mirzakhani's research are summarized in Table \ref{tab:mirzakhani_results}.

\begin{table}[htbp]
\centering
\caption{Summary of Mirzakhani's main contributions and methodologies.}
\label{tab:mirzakhani_results}
\vspace{0.2cm}
\small
\begin{tabular}{|l|l|}
\hline
\textbf{Result} & \textbf{Methodological Approach} \\
\hline
Simple geodesic counting & Polynomial growth $L^{6g-6}$ via WP volumes \\
\hline
WP volume recursion & Generalized McShane identity integration \\
\hline
Witten conjecture proof & Laplace transform and KdV hierarchy \\
\hline
Large genus asymptotics & Asymptotic analysis for $g \to \infty$ \\
\hline
Orbit closure classification & Affine invariant submanifolds (Magic Wand) \\
\hline
\end{tabular}
\end{table}


Mirzakhani’s work has left a transformative mark on modern mathematics, providing tools that continue to drive research in enumerative geometry and dynamical systems. From the perspective of \textit{Mathematics Education}, her legacy emphasizes the role of **visual intuition** as a rigorous heuristic tool. By reducing abstract problems on high-dimensional moduli spaces to the structural decomposition of surfaces, she demonstrated that geometric imagination is not merely illustrative, but a fundamental component of mathematical discovery. Her methodology serves as a testament to the synergy between visual creativity and formal rigor.